% Exposición académica de las dos narrativas conservadas en el cuaderno.
% Procedencia y correspondencia de contenidos: gestion/mapa_narrativa.md.
% H_lazo y rho_lazo no designan la holonomía nonádica Gamma_9.
\section{Compatibilidad estructural de la acción, la masa y la información operacional}
\label{sec:narrativa-estructural}

La composición de las ecuaciones permite identificar una relación más precisa
que la equivalencia entre masa y energía o la proporcionalidad entre energía
y frecuencia: una misma estructura de ruta determina una frecuencia propia,
una relación entre longitudes y una ponderación térmica. El par angular
participa en estas tres lecturas porque contiene el coeficiente de acción y
determina los canales de transporte. La exposición que sigue conserva
conjuntamente la construcción de los objetos, sus operaciones de lectura y
las condiciones de sus realizaciones físicas.

Esta composición permite, además, caracterizar la huella que deja el orden
de una secuencia de operaciones. Dicha huella admite una expresión mediante
el par angular y una aplicación inversa que restituye el coeficiente de
acción. En el ámbito constitutivo, la respuesta del vacío y el momento
espectral de Catalán se reúnen mediante una misma colección funcional. El
contenido explicativo de estas relaciones reside en las restricciones que
imponen conjuntamente, y no en la mera acumulación de identidades.

La genealogía APP--TRIT--TPK proporciona los objetos que se componen. Las
lecturas posteriores conservan aspectos distintos de una misma estructura:
escala, orientación, fase, memoria, incidencia y respuesta. Su análisis
conjunto establece qué relaciones permanecen determinadas al cambiar de
representación, qué información elimina una proyección y qué observaciones
complementarias permiten restituirla. Las demostraciones de las composiciones
y de sus dominios se desarrollan en las secciones demostrativas.

\subsection{La acción en la composición angular y en la estructura electrónica}
\label{subsec:narrativa-accion-angular}

La sección de Planck interviene antes de la evaluación energética final.
Su coeficiente participa en la geometría angular utilizada posteriormente
por la ley de masas. El enlace, con las unidades conservadas, es
\begin{equation}
 \hbar_{\mathrm{ret}}
   =\eta_{\mathrm{ret}}\,10^{-34}\mathcal U_S,
 \qquad
 A_{\deg}=1000\alpha,
 \qquad
 C^*_{\deg}=2(\eta_{\mathrm{ret}}+\alpha).
 \label{eq:narrativa-par-angular}
\end{equation}
El coeficiente $\eta_{\mathrm{ret}}$ es adimensional y por ello puede
componerse con $\alpha$ en la expresión angular. La constante dimensional
$\hbar_{\mathrm{ret}}$ conserva su unidad de acción. Esta distinción permite
seguir una misma dependencia sin sumar magnitudes de dimensiones diferentes.

La acción desempeña así dos funciones relacionadas. Por una parte, fija una
escala dimensional; por otra, participa en la estructura que determina el
reparto del transporte entre las dos orientaciones. Después de la conversión
única de grados a radianes, el ángulo y el contraángulo producen los canales
$q_+$ y $q_-$. Estos canales intervienen en las incidencias hexádico--octádicas,
en el funcional de Barbero--Immirzi, en la respuesta constitutiva y en los
términos de la ley de masas.

La sustitución de una constante por su expresión estructural hace visibles
dependencias que permanecen implícitas cuando la ecuación se escribe mediante
constantes aisladas. La expresión ampliada transmite la operación que
determina el valor, las coordenadas de las que depende y las condiciones
bajo las cuales puede reutilizarse. Su función explicativa procede de esta
información transmitida a la etapa siguiente.

En el electrón, la expresión completa conserva el exponente
$\varphi/\pi^2$, las correcciones construidas con $\alpha$ y $\Delta_4$,
y el factor de retorno
\begin{equation}
 \left(\frac{H_5}{\eta_{\mathrm{ret}}}\right)^{80-54\alpha+6\Delta_4}.
 \label{eq:narrativa-factor-electronico}
\end{equation}
El mismo $\eta_{\mathrm{ret}}$ que interviene en el contraángulo participa
en la lectura electrónica completa. La geometría angular y la corrección
de acción del electrón comparten, por tanto, una dependencia efectiva.
Estudiar conjuntamente ambas ecuaciones permite determinar qué relaciones
quedan impuestas por la intervención de una misma coordenada estructural
en operaciones distintas.

\subsection{Frecuencia propia y normalización estructural de la masa}
\label{subsec:narrativa-masa-frecuencia}

La ley general de masas conserva su prefactor electrónico. Se designa por
$\Xi_\Gamma$ el factor estructural completo de una ruta $\Gamma$: reúne la
normalización electrónica, el carácter de la ruta con su firma y sus torres,
y la potencia correspondiente de la razón de acción
$H_5/\eta_{\mathrm{ret}}$. Esta notación abrevia una composición de
operaciones ya definidas; su contenido permanece ligado a esas operaciones.

El reloj común tiene frecuencia angular
$\omega_0=2\pi/(108t_0)$, con $t_0>0$. La composición de la ley de masas con la sección
de acción da
\begin{equation}
 \boxed{
 m_\Gamma=\frac{\hbar_{\mathrm{ret}}\omega_0}{c^2}\,\Xi_\Gamma,
 \qquad
 E_\Gamma=\hbar_{\mathrm{ret}}\omega_0\Xi_\Gamma,
 \qquad
 \Omega_\Gamma=\frac{E_\Gamma}{\hbar_{\mathrm{ret}}}
             =\omega_0\Xi_\Gamma.}
 \label{eq:narrativa-masa-energia-frecuencia}
\end{equation}
Aquí $E_\Gamma$ es la energía de reposo y $\Omega_\Gamma$ su frecuencia
energética asociada.

Esta lectura separa el ritmo común de la contribución propia de cada
estructura. Dos secciones pueden compartir el período elemental y presentar
masas y frecuencias propias diferentes: sus firmas, retornos y memorias
producen factores $\Xi_\Gamma$ diferentes. La diversidad de las lecturas
específicas queda referida a la estructura de la sección, mientras el reloj
común establece su normalización temporal.

En el cociente $E_\Gamma/\hbar_{\mathrm{ret}}$ se cancela el factor
absoluto de acción que aparece explícitamente en
\eqref{eq:narrativa-masa-energia-frecuencia}. Permanece, sin embargo, su
participación estructural en los canales y en la razón de retorno contenida
en $\Xi_\Gamma$. La cancelación de un factor en una lectura posterior
debe distinguirse de la eliminación de las relaciones mediante las cuales
ese factor se construyó y volvió a intervenir.

Los cocientes entre masas coinciden con los cocientes entre estas frecuencias
propias. Una modificación común de las unidades de acción, tiempo o velocidad
deja inalterados tales cocientes. La estructura relativa tiene, por ello,
una exigencia de contraste propia: sus relaciones entre secciones deben
mantenerse coherentes, con independencia de la normalización común elegida.

La frecuencia asociada a una fase levantada conserva, además, el número de
vueltas de la historia. Dos fases iguales después de su reducción módulo
$2\pi$ pueden corresponder a levantamientos diferentes. El retorno de fase
conserva una regularidad, mientras la historia distingue el recorrido que
produjo ese retorno. La identificación del estado completo con la fase
reducida eliminaría precisamente esta distinción.

\subsection{Dilatación recíproca de las lecturas espaciales de la masa}
\label{subsec:narrativa-reciprocidad-espacial}

En la realización gravitatoria circular, la acción y la gravitación
satisfacen
\begin{equation}
 G=\frac{c^3L_\alpha^2}{\hbar_{\mathrm{ret}}},
 \qquad R_{108}=L_\alpha=\frac{c}{\omega_0}.
 \label{eq:narrativa-gravedad-accion}
\end{equation}
La composición con la masa de
\eqref{eq:narrativa-masa-energia-frecuencia} determina
\begin{equation}
 \boxed{
 r_{g,\Gamma}=R_{108}\Xi_\Gamma,
 \qquad
 \bar\lambda_\Gamma=\frac{R_{108}}{\Xi_\Gamma},
 \qquad
 r_{g,\Gamma}\bar\lambda_\Gamma=R_{108}^{2}.}
 \label{eq:narrativa-longitudes-reciprocas}
\end{equation}
Se utiliza la convención de radio gravitatorio reducido del desarrollo;
la construcción de $L_\alpha$ y la normalización conjunta están
demostradas en la sección~\ref{md:rigidez}. El reloj y el radio son
dos expresiones compatibles de esa misma longitud.
$\bar\lambda_\Gamma$ designa la longitud de Compton reducida. Estas
identidades recíprocas se formulan para $\Xi_\Gamma>0$; la lectura
$R_{108}/\Xi_\Gamma$ requiere ese dominio másico positivo.

El mismo factor que aumenta la frecuencia propia dilata la lectura
gravitatoria y contrae la longitud asociada a la acción y la masa. Las dos
transformaciones conservan un área común. En esta realización, el factor
másico actúa como parámetro de dilatación recíproca: una longitud crece con
él y la otra disminuye inversamente.

Esta relación proporciona un contenido matemático preciso a la vinculación
entre masa y configuración espacial. La jerarquía de la composición es
esencial: primero se determina la ruta y su factor estructural; después se
obtienen sus lecturas temporales, energéticas y espaciales. La estructura
se transmite a las lecturas mediante operaciones determinadas, en lugar de
reconstruirse retrospectivamente a partir de una masa medida.

La condición circular pertenece al resultado
\eqref{eq:narrativa-longitudes-reciprocas}. Las dos longitudes son lecturas
definidas por esa realización; su interpretación como radios materiales de
una partícula requeriría la identificación física correspondiente. La
igualdad establece aquí una reciprocidad exacta entre las magnitudes
definidas, con la coordenatización de acción, el reloj y la velocidad conservados.

La incorporación del momento mantiene esta masa estructural en la ecuación
de propagación. En la realización libre correspondiente, la frecuencia
propia determina el término másico de la dispersión, mientras el momento
describe la variación espacial. Así quedan diferenciadas la estructura
persistente de la sección y su estado de movimiento.

\subsection{Orientación de los canales y respuesta constitutiva del vacío}
\label{subsec:narrativa-vacio-orientacion}

La composición de permitividad, permeabilidad y funcional de Barbero--Immirzi
permite distinguir el valor de una respuesta de la estructura que lo
determina. El lector constitutivo utiliza dos canales etiquetados, con
respuestas $r_+$ y $r_-$. La velocidad normalizada depende de su producto;
la impedancia depende de su cociente. La velocidad conserva una relación
común entre los canales, mientras la impedancia distingue el reparto de la
respuesta entre sus orientaciones.

Dentro de la colección normalizada, la velocidad y el funcional de
Barbero--Immirzi permiten recuperar conjuntamente las dos respuestas
etiquetadas. El funcional aporta información orientada que la velocidad
aislada no determina. La restitución se apoya en la composición completa
de sus lectores y conserva la distinción entre cada canal y su agregado.

Sean $x$ e $y$ las coordenadas en radianes del ángulo y del contraángulo.
Para $x>0$ y $|y|<x$, la respuesta satisface
\begin{equation}
 \widehat c(x,y)\geq\widehat c(x,0),
 \label{eq:narrativa-convexidad-velocidad}
\end{equation}
con igualdad únicamente cuando $y=0$. Más precisamente, manteniendo $x$
fijo, $\log\widehat c$ es una función par y estrictamente convexa de $y$.
La demostración de esta propiedad pertenece al desarrollo constitutivo
posterior.

La velocidad pierde el signo de la orientación, pero conserva un efecto
de la separación entre los canales. Cerca de la configuración simétrica,
ese efecto aparece a segundo orden; la impedancia distingue la orientación
ya a primer orden. Por ello, sustituir ambos canales por su media antes de
evaluar la respuesta cambia el resultado. Una lectura escalar puede ocultar
una orientación y conservar, al mismo tiempo, una consecuencia de ella.

Este comportamiento especifica qué información transmite cada operación.
La pérdida del signo, la conservación de una separación y la recuperación
del reparto constituyen propiedades diferentes. La estructura mantiene
consecuencias observables incluso cuando una de sus lecturas se expresa
mediante un único número.

La desigualdad \eqref{eq:narrativa-convexidad-velocidad} describe la colección
matemática de respuesta en su dominio. Una evolución física debe situarse
en los estados y trayectorias efectivamente producidos por TPK. La
parametrización analítica permite estudiar la sensibilidad del lector;
la identificación de una variación paramétrica con una trayectoria física
conserva su obligación de realización.

\subsection{Boltzmann, frecuencia estructural y ponderación térmica ternaria}
\label{subsec:narrativa-boltzmann}

La sección térmica relaciona la energía del ciclo con la constante de
Boltzmann y con su temperatura correspondiente:
\begin{equation}
 k_B\Theta_{108}\log3=E_{\mathrm{clk}}.
 \label{eq:narrativa-boltzmann-ciclo}
\end{equation}
Como la energía de una sección es
$E_\Gamma=E_{\mathrm{clk}}\Xi_\Gamma$, el peso térmico evaluado a esa
temperatura resulta
\begin{equation}
 \boxed{\exp\!\left(-\frac{E_\Gamma}{k_B\Theta_{108}}\right)
       =3^{-\Xi_\Gamma}.}
 \label{eq:narrativa-peso-termico}
\end{equation}
La base ternaria aparece por composición de la normalización informacional
con la acción y el reloj. La estructura de ruta que determina la masa
determina también el exponente del peso térmico.

El peso y la probabilidad desempeñan funciones distintas. Las probabilidades
se obtienen al normalizar los pesos y conservar las multiplicidades de los
estados. Asimismo, varias historias equivalentes requieren una identificación
de estados antes de recibir una multiplicidad física: el número de historias
y la degeneración del estado corresponden a objetos diferentes.

La composición permite reunir las lecturas en una identidad común:
\begin{equation}
 \boxed{
 \Xi_\Gamma
 =\frac{\Omega_\Gamma}{\omega_0}
 =\frac{r_{g,\Gamma}}{R_{108}}
 =\frac{R_{108}}{\bar\lambda_\Gamma}
 =\frac{E_\Gamma}{k_B\Theta_{108}\log3}.}
 \label{eq:narrativa-cuatro-lecturas}
\end{equation}
Se obtienen cuatro formas coordinadas de leer una misma componente
estructural: temporal, espacial, gravitatoria y térmica. Cada igualdad
conserva las condiciones de su realización, especialmente la condición
circular gravitatoria y la temperatura del ciclo.

La identidad \eqref{eq:narrativa-cuatro-lecturas} especifica la relación
que comparten las magnitudes y proporciona un criterio para contrastar su
compatibilidad. La masa relativa puede caracterizarse mediante la frecuencia
relativa; la misma coordenada determina una dilatación recíproca de
longitudes y una ponderación térmica. Las funciones de lectura son distintas,
pero su dependencia queda vinculada por la genealogía común.

La relación térmica identifica el estado y la normalización en los que
interviene $\Theta_{108}$. Esta identificación mantiene separadas la energía
de la sección y su eventual transferencia como calor; también mantiene
localizado el significado de la temperatura del ciclo, sin extenderla al
universo como temperatura única. La capacidad explicativa procede de las
condiciones explícitas de la composición.

\subsection{Elipse de acción y transporte de la forma geométrica}
\label{subsec:narrativa-elipse-transporte}

La elipse de acción tiene semiejes cuyo producto es $2\hbar$. El área
orientada acumulada durante una vuelta es $h=2\pi\hbar$. Una deformación
elíptica redistribuye los semiejes conservando su producto. En la
realización eléctrica y magnética se obtienen dos contribuciones
proporcionales a $\cos^2\theta$ y $\sin^2\theta$, cuya suma es
$\hbar\omega$. La estructura determina, por tanto, el reparto de la
energía durante el ciclo, además de su valor total.

El transporte de una elipse cuya forma cambia exige conservar la relación
entre la deformación y la acción. Sea $\zeta$ la coordenada de deformación,
determinada por la razón entre el contraángulo y el ángulo. Para una
deformación diferenciable, el transporte exacto introduce
\begin{equation}
 H_{\mathrm{tot}}
 =\frac{\omega}{2}
   \left(e^{-\zeta}Q^2+e^\zeta P^2\right)
  +\frac{\dot\zeta}{2}QP.
 \label{eq:narrativa-hamiltoniano-transporte}
\end{equation}
El primer término describe la evolución en una elipse de forma fija.
El segundo procede del transporte de una forma variable. La demostración
de la identidad de transporte establece que este término compensa
exactamente la variación geométrica y conserva la acción cuadrática
correspondiente. Su omisión o la inversión de su signo destruyen esa
compensación.

La interpretación exige distinguir dos situaciones. Cuando se cambia
únicamente de coordenadas, el término adicional es una conexión de
representación. Un intercambio físico requiere, en cambio, una variación
relativa entre sistema y referencia. La reexpresión de un Hamiltoniano
conserva el contenido de la representación; la atribución de una interacción
necesita determinar las operaciones relativas que actúan sobre el sistema.

El par angular rector conserva otra dependencia decisiva: fijar $\alpha$
y la sección de acción fija también $\zeta$. Una evolución de la forma
debe, por ello, extraerse de los estados y operaciones efectivos. La
deformación paramétrica por sí sola describe una colección de representaciones;
su identificación como evolución HMT transmite las relaciones que producen
el ángulo y el contraángulo.

El transporte conjunto de la forma y de la sección de acción conserva la
acción normalizada, mientras el área sin normalizar cambia de acuerdo con
el cociente entre las dos secciones. Esta construcción distingue tres
operaciones: redistribuir una acción conservada, transportar una sección de
acción y cambiar la representación o las unidades. La distinción permite
examinar interacciones conservando un balance que incluya el origen de la
variación, en lugar de atribuir generación de energía a una mera deformación
de coordenadas.

\subsection{Compatibilidad conjunta y contraste relacional}
\label{subsec:narrativa-compatibilidad-conjunta}

El eje de investigación es la determinación de las condiciones de
compatibilidad que una misma estructura impone a sus diferentes
manifestaciones físicas. Dos líneas de composición se complementan.

La primera reúne relaciones estáticas: los observables complementarios
restituyen los canales, y la estructura recuperada determina la acción,
la respuesta constitutiva y los factores de masa. El contraste pertinente
es conjunto. Una modificación que preserve una cifra debe conservar también
las demás relaciones que utiliza la misma construcción. La velocidad,
Barbero--Immirzi y la impedancia proporcionan un sistema de comprobación;
las razones de masas y de frecuencias proporcionan otro.

La segunda línea estudia las historias efectivas y compara sus resultados
después de transportar coherentemente el sistema y su referencia. Las
diferencias de fase relativa, respuesta constitutiva relativa o transferencia
entre modos tienen un contenido distinto de las diferencias debidas al
nombre de las coordenadas. La memoria, la orientación y los retornos
intervienen mediante las operaciones que una lectura escalar puede eliminar.

Una relación estructural alcanza capacidad explicativa y de contraste
cuando varias lecturas independientes deben concordar por proceder de una
misma estructura. Su valor reside en la rigidez de esa concordancia:
acción, ritmo, respuesta espacial y ponderación informacional dejan de
poder elegirse por separado dentro de la realización declarada. La
composición reúne esas restricciones y permite demostrar sus consecuencias
sin sustituir su genealogía por ajustes individuales de los resultados.

\subsection{Cierre de fase y holonomía relativa del estado}
\label{subsec:narrativa-holonomia-relativa}

El ciclo dodecafásico contiene el plano en el que se realiza la raíz
interna de $-1$. La inclusión de ese plano en las doce componentes es
explícita; el cuarto de giro utilizado a continuación procede de dicha
construcción. Sobre él actúa la forma de la elipse de acción determinada
por el par angular, cuya coordenada es
\begin{equation}
 \xi=\frac{C^*}{A},\qquad
 \zeta=\operatorname{artanh}\xi.
 \label{eq:narrativa-coordenada-forma}
\end{equation}
El cociente utiliza los dos ángulos expresados en la misma unidad y tiene
el mismo valor en grados y en radianes. La coordenada real de forma se
define para $|\xi|<1$. La forma conjugada corresponde al
cambio de signo de $\zeta$.

Se compone el cuarto de giro de una forma con el de su conjugada y con
sus inversos. Si $U_+$ y $U_-$ designan las dos operaciones, la
composición exacta es
\begin{equation}
 \boxed{
 H_{\mathrm{lazo}}=U_+U_-U_+^{-1}U_-^{-1}
 =\begin{pmatrix}\rho_{\mathrm{lazo}}&0\\[2pt]
                  0&\rho_{\mathrm{lazo}}^{-1}\end{pmatrix},
 \qquad
 \rho_{\mathrm{lazo}}=
 \left(\frac{1+\xi}{1-\xi}\right)^2.}
 \label{eq:narrativa-holonomia-lazo}
\end{equation}
La composición dilata una coordenada y contrae la conjugada por el factor
inverso: conserva el área y modifica el estado. Para $\xi\ne0$, su
resultado difiere de la identidad.

Las cuatro operaciones incluyen sus inversas y sus incrementos de fase
declarados suman cero. El orden de composición conserva, sin embargo, un
efecto. Una descripción que sólo retenga los incrementos que se añaden y
se compensan elimina ese efecto relativo. El cierre de la contabilidad
de fase y el retorno del estado completo son, por tanto, propiedades
distintas de la secuencia.

La igualdad admite un levantamiento al espacio original de doce
componentes. El operador actúa mediante la dilatación recíproca sobre el
plano considerado y como identidad sobre su complemento. El enlace con
el ciclo dodecafásico se realiza mediante una transformación explícita;
el plano conserva así su procedencia dentro del espacio de fases.

La notación $H_{\mathrm{lazo}}$ designa esta composición posterior y
mantiene diferenciada la holonomía nonádica $\Gamma_9$. La realización
ilustra de manera concreta una distinción estructural: el cierre de una
lectura del recorrido puede coexistir con una transformación persistente
del estado. La procedencia de los operadores y el alcance de su
levantamiento permanecen expuestos en el desarrollo demostrativo.

\subsection{Restitución del coeficiente de acción desde el transporte relativo}
\label{subsec:narrativa-restitucion-accion}

El par angular de \eqref{eq:narrativa-par-angular} puede expresarse con
$s_0=10^{-34}\mathcal U_S$ y
$\hbar_{\mathrm{ret}}=s_0\eta_{\mathrm{ret}}$. El coeficiente
$\eta_{\mathrm{ret}}$ conserva su carácter adimensional, mientras $s_0$
conserva la escala dimensional. Al incorporar el contraángulo completo
en \eqref{eq:narrativa-holonomia-lazo}, se obtiene
\begin{equation}
 \boxed{
 \rho_{\mathrm{lazo}}=
 \left(\frac{501\alpha+\eta_{\mathrm{ret}}}
              {499\alpha-\eta_{\mathrm{ret}}}\right)^2.}
 \label{eq:narrativa-lazo-alfa-accion}
\end{equation}
Esta expresión utiliza la rama positiva
$0<\eta_{\mathrm{ret}}<499\alpha$. Los números $501$ y $499$ proceden
de $500\pm1$, y $500$ procede de $1000/2$: son consecuencias de las
normalizaciones que ya intervienen en el par angular.

La relación entre $\alpha$ y el coeficiente de acción determina la
separación entre las respuestas conjugadas después del recorrido
compuesto. La acción participa tanto en la escala dimensional como en
la forma que gobierna el transporte. La fórmula hace explícita esta
doble intervención sin identificar ambas funciones.

La relación admite la inversión
\begin{equation}
 \boxed{
 \eta_{\mathrm{ret}}=
 \alpha\left[
 500\,\frac{\sqrt{\rho_{\mathrm{lazo}}}-1}
               {\sqrt{\rho_{\mathrm{lazo}}}+1}-1\right].}
 \label{eq:narrativa-inversa-accion}
\end{equation}
Con $\alpha$ y los ejes orientados conservados, la respuesta del lazo
permite restituir el coeficiente de acción. La lectura es estrictamente
creciente respecto de $\rho_{\mathrm{lazo}}$ en la rama declarada.

Se obtiene una caracterización operacional: el coeficiente de acción
puede restituirse mediante la dilatación relativa producida por las
formas conjugadas del par angular. Su definición generadora permanece
en la construcción anterior; la caracterización operacional establece
cómo recuperarlo en una representación distinta.

La escala $s_0$ sigue siendo necesaria para recuperar la constante
dimensional $\hbar_{\mathrm{ret}}$. La forma del transporte conserva
una relación adimensional; la escala de acción determina la
normalización física completa. Esta separación identifica con precisión
qué dato devuelve la inversión y qué estructura adicional transmite la
lectura dimensional.

\subsection{Iteración de la holonomía e información operacional de la historia}
\label{subsec:narrativa-informacion-operacional}

La repetición del lazo conserva el mismo programa de operaciones y
produce una ley explícita de acumulación:
\begin{equation}
 H_{\mathrm{lazo}}^n=
 \begin{pmatrix}\rho_{\mathrm{lazo}}^n&0\\
                 0&\rho_{\mathrm{lazo}}^{-n}\end{pmatrix}.
 \label{eq:narrativa-iteracion-lazo}
\end{equation}
Una preparación $(Q_0,P_0)$ se transforma en
\begin{equation}
 Q_n=\rho_{\mathrm{lazo}}^nQ_0,
 \qquad
 P_n=\rho_{\mathrm{lazo}}^{-n}P_0.
 \label{eq:narrativa-iteracion-estado}
\end{equation}
El producto $Q_nP_n$ permanece constante y el transporte conserva
el área simpléctica, mientras cambia la proporción entre las dos
coordenadas.

Una regla finita repetida puede dejar una huella acumulativa que
distinga sus iteraciones. Una preparación conocida y una referencia
conservada permiten recuperar $n$ mediante el incremento logarítmico
de una coordenada no nula, para $\rho_{\mathrm{lazo}}\ne1$.
Cuando se conserva
únicamente el estado final y se desconoce la preparación, la misma
restitución queda indeterminada.

La información accesible depende de la lectura retenida: balance,
amplitud, proporción orientada, preparación u operador completo. Cada
elección puede identificar historias que otra lectura distingue.
La comparación de
\begin{equation}
 U_+U_-U_+^{-1}U_-^{-1}
 \qquad\text{y}\qquad
 U_+U_+^{-1}U_-U_-^{-1}
 \label{eq:narrativa-orden-multiplicidades}
\end{equation}
proporciona un contraste directo. Ambos productos contienen las mismas
operaciones con las mismas multiplicidades. El segundo da la identidad;
el primero da $H_{\mathrm{lazo}}$. Una estadística de frecuencias de
aparición resulta insuficiente para predecir esta respuesta: debe
conservarse también el orden.

Este resultado precisa qué descripción estadística es insuficiente
para el objeto estudiado y qué observaciones restituyen la diferencia.
La teoría de información de Shannon admite distribuciones sobre
secuencias ordenadas y puede representar ese orden. La distinción
operacional se refiere al contenido retenido por el observable elegido.

La información operacional de una historia se caracteriza mediante
su capacidad de producir respuestas distinguibles para preparaciones
y detectores especificados. Esta definición permite precisar cuánto
conserva una lectura. Su alcance corresponde a la representación
utilizada: historias distintas pueden inducir un mismo operador
reducido y conservar diferencias en otras componentes de la genealogía
TPK.

La traza de $H_{\mathrm{lazo}}$ coincide con la de
$H_{\mathrm{lazo}}^{-1}$. Detecta la intensidad del efecto, pero elimina
su orientación. Una lectura que conserve los ejes y compare sus
respuestas distingue los dos sentidos. La elección de detector forma
parte, por tanto, de la restitución del recorrido.

La finitud del programa, el crecimiento de una memoria observable y
una tasa entrópica describen propiedades diferentes. Una tasa de
crecimiento combinatorio nula puede coexistir con estados distinguibles.
La entropía de la dinámica completa requiere especificar su espacio
de estados y, cuando corresponda, su medida. La composición demuestra
la acumulación y su restitución en la representación declarada;
la evaluación de una entropía termodinámica conserva sus variables
y condiciones físicas propias.

\subsection{Potencial espectral y reciprocidad diferencial de las respuestas}
\label{subsec:narrativa-potencial-espectral}

Los canales constitutivos conservan las coordenadas
$x=A_{\mathrm{rad}}$, $y=C^*_{\mathrm{rad}}$ y las incidencias $90/120$.
Su composición determina la identidad
\begin{equation}
 \boxed{
 \frac{\partial\log\widehat c}{\partial y}
 =\frac{\partial\log\widehat Z}{\partial x}.}
 \label{eq:narrativa-reciprocidad-diferencial}
\end{equation}
Las magnitudes $\widehat c$ y $\widehat Z$ designan la velocidad y la
impedancia normalizadas del lector.

La sensibilidad de la propagación a la separación orientada queda
ligada a la sensibilidad de la impedancia a la coordenada común.
Las dos respuestas comparten una función espectral y sus variaciones
deben mantener esa relación para seguir perteneciendo a la misma
construcción.

Existe una función $\mathscr F(x,y)$ cuyas derivadas son
\begin{equation}
 \partial_x\mathscr F=\log\widehat c,
 \qquad
 \partial_y\mathscr F=\log\widehat Z.
 \label{eq:narrativa-gradiente-potencial}
\end{equation}
Se trata de un potencial espectral definido desde los canales $90/120$,
con serie convergente y normalización explícita. Su Hessiano es
negativo definido en $x>|y|$, propiedad que permite demostrar
unicidad y sensibilidad de la restitución.

La igualdad entre las derivadas cruzadas establece una restricción
sobre la variación conjunta de las respuestas dentro de esta colección
de realización. Una modificación que destruya la identidad
\eqref{eq:narrativa-reciprocidad-diferencial} cambia el lector
constitutivo. El contraste afecta así a las relaciones de variación,
además de a los valores evaluados.

El potencial es de valor único: la integral de su diferencial sobre
un lazo cerrado en las dos coordenadas se anula. La composición
ordenada de transportes puede conservar, en cambio, una holonomía
no trivial. La proyección constitutiva y el transporte operatorio
retienen aspectos diferentes de la historia; su comparación exige
conservar sus dominios y operaciones de lectura.

\subsection{Momento de Catalán y respuesta contractiva de las incidencias}
\label{subsec:narrativa-catalan-potencial}

La conexión entre el momento orientado de Catalán y el potencial
constitutivo se expresa mediante el momento de profundidad dos
\begin{equation}
 D_2(z)=\sum_{n\geq1}\frac{z^n}{n^2}.
 \label{eq:narrativa-momento-dos}
\end{equation}
El momento de Catalán se obtiene como la parte imaginaria de
$D_2(is)$; en $s=1$ recupera la constante de Catalán. El potencial
constitutivo utiliza la misma función sobre las contracciones de
los canales, en $q_\pm^{90}$ y $q_\pm^{120}$, con las normalizaciones
correspondientes.

El cuarto de giro orientado y la respuesta contractiva $90/120$
se reúnen mediante un mismo momento espectral, evaluado sobre
soportes distintos. Catalán adquiere aquí significado por la
operación espectral que representa y por la relación de esa
operación con la respuesta constitutiva. La colección funcional
conserva información que excede la de su valor extremo.

La composición mantiene también las diferencias entre sus objetos.
El valor aislado de Catalán no determina todos los canales;
la relación utiliza el momento espectral completo, sus argumentos
y la orientación. La convergencia de la serie, las identidades
de sus derivadas y su compatibilidad con la amplificación forman
parte del desarrollo demostrativo de esta conexión.

\subsection{Restitución exacta y estabilidad ante precisión finita}
\label{subsec:narrativa-estabilidad-restitucion}

El par $(\widehat c,\widehat Z)$ recupera de manera única los dos
parámetros del lector en el dominio declarado. Las cotas de
sensibilidad relacionan, además, el error de las respuestas con
el error de los parámetros restituidos. La unicidad exacta y la
estabilidad de la inversión constituyen resultados distintos y
complementarios.

Esta distinción separa una información eliminada por la proyección
de una información conservada pero difícil de resolver con precisión
finita. En regiones de contracciones intensas, variaciones apreciables
de los parámetros pueden producir diferencias muy pequeñas en las
respuestas. La inversión sigue siendo única para datos exactos,
mientras una incertidumbre finita puede impedir distinguir tales
variaciones.

La restitución genealógica exige especificar tanto qué estructura
puede recuperarse como la estabilidad con la que puede recuperarse.
El margen de restitución designa la cota que cuantifica esa estabilidad
en un dominio. La elección de observable debe atender a la diferencia
que se pretende resolver: una lectura puede ser correcta y resultar
poco sensible a ella.

La traza del lazo ofrece un ejemplo de esta selección. Su lectura
elimina el sentido del recorrido; la respuesta sobre ejes orientados
lo conserva. La determinación del detector adecuado pertenece a
la demostración de recuperabilidad y al diseño del contraste.

\subsection{Síntesis: coordenadas de la sección y memoria de las operaciones}
\label{subsec:narrativa-sintesis}

La identidad \eqref{eq:narrativa-cuatro-lecturas} conserva como eje
la relación entre el factor estructural de una sección, su frecuencia
propia, las dos longitudes recíprocas y su lectura térmica. Permanecen
explícitas las condiciones: la misma coordenatización de acción, la realización
gravitatoria circular y la temperatura del ciclo.

La composición operatoria añade una distinción central. El factor
de la sección determina sus lecturas coordinadas; el transporte
determina el resultado de componer operaciones sobre ella y la
huella que conserva su orden. El factor másico y la dilatación
$\rho_{\mathrm{lazo}}$ son objetos diferentes. Su relación debe
pasar por los lectores correspondientes para distinguir una
transformación del estado de una identificación másica.

Planck interviene en la normalización de la acción y en la forma
angular del transporte. Boltzmann interviene en la relación entre
energía y ponderación térmica. La reciprocidad gravitatoria conserva
el producto de las longitudes en su realización declarada. Velocidad
e impedancia proporcionan lecturas complementarias de los canales.
Esta arquitectura admite conjuntamente una descripción de
compatibilidades y una descripción de la memoria operacional.

La capacidad explicativa de la estructura reside en determinar
varias respuestas, imponer relaciones entre ellas y especificar
qué historia puede restituirse desde sus observaciones. Las
coincidencias de valores quedan acompañadas por operaciones,
condiciones de compatibilidad y criterios de restitución.

La composición distingue el cambio pasivo de coordenadas, que
telescopa cuando retornan coordenatización y fase, del transporte relativo
cuyo operador final es $H_{\mathrm{lazo}}\ne I$. Un intercambio
físico requiere realizar las operaciones, sus inversas y la
referencia en un mismo sistema, incluyendo el controlador en
el balance. Este criterio conserva el contenido matemático
de una variación energética sin atribuirle producción de energía
por el mero cambio de lectura.

Las pruebas de composición, iteración, restitución del coeficiente
de acción, reciprocidad constitutiva y sensibilidad sostienen
estas caracterizaciones en sus dominios respectivos. Narración,
definiciones y demostraciones forman así una exposición continua:
la primera identifica el significado de las relaciones, las
segundas fijan los objetos y las terceras establecen las
operaciones que permiten componerlos y recuperarlos.
